% =========================================================
% Boolean Rings
% =========================================================
\paragraph{Boolean Rings}

\begin{definitionbox}{Definition --- Boolean Ring}
\begin{definition}[Boolean Ring]
\label{def:boolean-ring}
A \textbf{Boolean ring} is a ring with unity $(R,+,\cdot,1)$ such that
every element is idempotent under multiplication:
\[
  \tag{B1}
  p\cdot p = p \quad \text{for every } p\in R.
\]
\end{definition}
\end{definitionbox}

\begin{remark*}[Standard quantified statement]
\[
\operatorname{BooleanRing}(R,+,\cdot,1)
\Longleftrightarrow
\operatorname{RingWithUnity}(R,+,\cdot,1)
\land
\forall p\in R,\; p\cdot p=p.
\]
\end{remark*}

\begin{remark*}[Definition predicate reading]
$(R,+,\cdot,1)$ is a Boolean ring exactly when it is a ring with a
multiplicative identity whose multiplication is idempotent.
\end{remark*}

\begin{remark*}[Negated quantified statement]
\[
\neg\operatorname{BooleanRing}(R,+,\cdot,1)
\Longleftrightarrow
\neg\operatorname{RingWithUnity}(R,+,\cdot,1)
\;\lor\;
\exists p\in R,\; p\cdot p\neq p.
\]
\end{remark*}

\begin{remark*}[Rewire note]
This definition reuses
\hyperref[def:idempotent-operation]{Idempotent Operation}:
\[
  \operatorname{BooleanRing}(R,+,\cdot,1)
  \Longleftrightarrow
  \operatorname{RingWithUnity}(R,+,\cdot,1)
  \land
  \operatorname{IdempotentOperation}(\cdot,R).
\]
There is no separate idempotent-element predicate here; the Boolean
condition is idempotence of multiplication as an operation.
\end{remark*}

\begin{remark*}[Characteristic 2 is derived, not assumed]
Some presentations list $p+p=0$ as an axiom of Boolean rings. In this
volume it is a theorem, Proposition~\ref{prop:boolean-ring-self-additive-inverse}.
Thus
\[
  \operatorname{BooleanRing}(R,+,\cdot,1)
  \Longrightarrow
  \operatorname{Characteristic2}(R,+).
\]
\end{remark*}

\begin{proposition}[Self Additive Inverse / Characteristic 2]
\label{prop:boolean-ring-self-additive-inverse}
In any Boolean ring $R$,
\[
  p+p=0 \quad \text{for every } p\in R.
\]
\hyperref[prf:boolean-ring-self-additive-inverse]{Go to proof.}
\end{proposition}

\begin{remark*}
\FlashQ{What is the additive characteristic of a Boolean ring?}
\FlashA{Every element satisfies $p+p=0$; Boolean rings have characteristic $2$ unless they are trivial.}
\FlashHint{Apply idempotence to $(p+q)^2$ and cancel the shared $p$ and $q$ terms.}
\end{remark*}

\begin{proposition}[Self Negation]
\label{prop:boolean-ring-self-negation}
In any Boolean ring $R$,
\[
  -p=p \quad \text{for every } p\in R.
\]
\hyperref[prf:boolean-ring-self-negation]{Go to proof.}
\end{proposition}

\begin{proposition}[Commutativity from Idempotence]
\label{prop:boolean-ring-commutativity-from-idempotence}
In any Boolean ring $R$,
\[
  pq=qp \quad \text{for all } p,q\in R.
\]
Hence every Boolean ring is a commutative ring
(Definition~\ref{def:notes-rings-definition-l65}).
\hyperref[prf:boolean-ring-commutativity-from-idempotence]{Go to proof.}
\end{proposition}

\begin{remark}[Structural Link to Integral Domains]
\label{rem:boolean-rings-and-integral-domains}
Every Boolean ring with more than two elements has zero divisors, hence
is not an integral domain (Definition~\ref{def:notes-rings-integral-domains-l25}).
Indeed, idempotence gives
\[
  p(1-p)=p-p^2=p-p=0.
\]
If $p\neq 0$ and $p\neq 1$, then $p$ and $1-p$ are nonzero zero divisors.
The only Boolean ring that is an integral domain is the two-element ring
$\mathbf{2}\cong\mathbb{F}_2$.
\end{remark}

\begin{example}[Standard Boolean Rings]
\label{ex:boolean-ring-standard-examples}
\begin{enumerate}[label=(\roman*)]
  \item The two-element ring
        \[
          \mathbf{2}=\{0,1\}
        \]
        with addition and multiplication modulo $2$ is a Boolean ring.
        Its operation tables are
        \[
        \begin{array}{c|cc}
        + & 0 & 1\\ \hline
        0 & 0 & 1\\
        1 & 1 & 0
        \end{array}
        \qquad
        \begin{array}{c|cc}
        \cdot & 0 & 1\\ \hline
        0 & 0 & 0\\
        1 & 0 & 1
        \end{array}
        \]
        and both elements satisfy $p^2=p$.

  \item For any $n\geq 1$, the product ring $\mathbf{2}^n$ with
        coordinatewise addition and multiplication is a Boolean ring.
        The four-element case $\mathbf{2}^2$ has elements
        $(0,0),(1,0),(0,1),(1,1)$, and idempotence holds coordinatewise:
        \[
          (a,b)^2=(a^2,b^2)=(a,b).
        \]

  \item For any set $X$, the power set
        \[
          \bigl(\mathcal P(X),\triangle,\cap,X\bigr)
        \]
        is a Boolean ring. Addition is symmetric difference, multiplication
        is intersection, the zero element is $\varnothing$, and the unity is
        $X$. This example is isomorphic to the function ring $\mathbf{2}^{X}$.
\end{enumerate}
\end{example}

\begin{theorem}[Stone Representation for Boolean Rings, stated]
\label{thm:stone-representation-boolean-rings}
Every Boolean ring is isomorphic to a field of sets. Every finite Boolean
ring is isomorphic to $\mathcal P(X)$ for some finite set $X$, hence has
$2^{|X|}$ elements.
\hyperref[prf:stone-representation-boolean-rings]{Go to proof.}
\end{theorem}

\begin{corollary}[Order is a Power of Two]
\label{cor:boolean-ring-finite-order-power-two}
A finite Boolean ring has order $2^n$ for some $n\geq 0$. In particular,
there is no Boolean ring of order $3$.
\hyperref[prf:boolean-ring-finite-order-power-two]{Go to proof.}
\end{corollary}

\begin{remark*}[Placement note]
The representation theorem is recorded here so the Boolean-ring facts stay
contiguous. A later abstract-algebra chapter may move the proof once ideals,
prime ideals, and ultrafilters are available.
\end{remark*}
